\section{Low-Cost Seller Profit}
\label{sec:low-cost-profit}

It suffices to focus on the low-type seller's profit. Suppose it sets reserve
price $\underline{r}$.

\subsection{When the auction fails}
\label{subsec:profit-fail}

The English auction \emph{fails} when every bidder stops before the reserve
$\underline{r}$ is met, so the sale moves to stage~2.
On the low bidding branch, failure is equivalent to $v_{\max}\le z$ with
\begin{equation}
  z
  \equiv \beta^{-1}(\underline{r})
  = \frac{\underline{r}-\delta\psi c_L}{1-\delta\psi}.
  \label{eq:z-def}
\end{equation}
This event has probability $\Pr(\text{fail})=F(z)^{N}$.

The seller may still earn positive profit in stage~2 if the selected
buyer makes a $c_H$ offer; each accepted $c_H$ offer yields profit $c_H-c_L$.
When the auction fails before $\underline{r}$, this event reveals no information
about the seller's type, so the posterior belief that the seller is low-cost
remains $\psi$.
By \eqref{eq:second-stage-threshold} with this posterior, a buyer with value $v$
offers $c_H$ in stage~2 if and only if $v\ge v_{\psi}$, where
\begin{equation}
  v_{\psi}
  \equiv \frac{c_H-\psi c_L}{1-\psi}.
  \label{eq:v-psi-def}
\end{equation}
The subscript records that here $\tilde{\psi}=\psi$.
Since failure implies $v_{\max}\le z$, the seller earns this margin only when
\[
  v_{\psi} < v_{\max} < z
  \qquad\text{(equivalently } z > v_{\max} > v_{\psi}\text{)}.
\]
Let $\underline{r}_0$ denote the reserve at which $z=v_{\psi}$,
\begin{equation}
  \underline{r}_0
  \equiv
  \delta\psi c_L
  + (1-\delta\psi)\frac{c_H-\psi c_L}{1-\psi}
  =
  \frac{(1-\delta\psi)c_H-\psi(1-\delta)c_L}{1-\psi}.
  \label{eq:failure-positive-cutoff}
\end{equation}
Since $z$ is increasing in $\underline{r}$, $v_{\psi}<v_{\max}<z$ is possible if
and only if $\underline{r}>\underline{r}_0$; otherwise $\Pi_{\text{fail}}=0$.\footnote{%
  The cutoff $\underline{r}_0$ need not lie below the top of the value support.
  If buyers place sufficiently high prior probability on the low-cost type,
  $v_{\psi}$ is so large that even the largest feasible $v_{\max}$ on a failure
  path need not warrant a $c_H$ offer.
  Indeed,
  $\underline{r}_0<1$ if and only if
  $\psi<(1-c_H)/\bigl(1-c_L-\delta(c_H-c_L)\bigr)$;
  if not, then $\underline{r}_0\ge1$ and $\Pi_{\text{fail}}=0$ for every feasible
  reserve $\underline{r}\le1$.}

For $\underline{r}>\underline{r}_0$, one bidder is drawn to make a
take-it-or-leave-it offer with probability $\delta$.
Here we distinguish between the highest bidder and the others, because stage~1
bidding reveals $v_{\max}$, which the seller observes, but not the other bidders'
values.

If the selected bidder is the highest, his value
equals $v_{\max}$, which the seller already knows, and if $v_{\max}\ge v_{\psi}$
he offers $c_H$.
This contributes
\begin{equation}
  \Pi_{\text{fail}}^{(1)}
  = \delta (c_H - c_L)
    \int_{v_{\psi}}^{z} dF(v)^{N}.
  \label{eq:profit-fail-max}
\end{equation}
If the selected bidder is not the highest, there are $N-1$ symmetric
possibilities.
The seller still knows $v_{\max}$ but not the selected value $v$, and a non-max
bidder with value $v$ offers $c_H$ with probability $1-F(v_{\psi})/F(v)$.
This contributes
\begin{equation}
  \Pi_{\text{fail}}^{(2)}
  = \delta (N-1) (c_H - c_L)
    \int_{v_{\psi}}^{z}
      \left(1 - \frac{F(v_{\psi})}{F(v)}\right)\, dF(v)^{N}.
  \label{eq:profit-fail-nonmax}
\end{equation}
Adding \eqref{eq:profit-fail-max} and \eqref{eq:profit-fail-nonmax}, and factoring
out $\delta(c_H-c_L)$,
\begin{equation}
  \Pi_{\text{fail}}
  = \Pi_{\text{fail}}^{(1)} + \Pi_{\text{fail}}^{(2)}
  = \delta (c_H - c_L)
    \Biggl[
      \int_{v_{\psi}}^{z} dF(v)^{N}
      + (N-1)
      \int_{v_{\psi}}^{z}
        \left(1 - \frac{F(v_{\psi})}{F(v)}\right)\, dF(v)^{N}
    \Biggr].
  \label{eq:profit-fail}
\end{equation}
The first integral in \eqref{eq:profit-fail} equals
$F(z)^{N}-F(v_{\psi})^{N}$.
For $N\ge 2$, integration by parts on the second integral gives
\begin{align}
  \int_{v_{\psi}}^{z}
    \left(1 - \frac{F(v_{\psi})}{F(v)}\right)\, dF(v)^{N}
  &= \left(1 - \frac{F(v_{\psi})}{F(z)}\right) F(z)^{N}
    - \frac{F(v_{\psi})}{N-1}
      \bigl(F(z)^{N-1} - F(v_{\psi})^{N-1}\bigr),
  \label{eq:profit-fail-nonmax-int}
\end{align}
since $\int_{v_{\psi}}^{z} F(v)^{N-2} f(v)\, dv
=\bigl(F(z)^{N-1}-F(v_{\psi})^{N-1}\bigr)/(N-1)$.
Adding the two integrals in \eqref{eq:profit-fail},
\begin{align}
  &\int_{v_{\psi}}^{z} dF(v)^{N}
    + (N-1)
    \int_{v_{\psi}}^{z}
      \left(1 - \frac{F(v_{\psi})}{F(v)}\right)\, dF(v)^{N}
  \notag\\
  &= F(z)^{N} - F(v_{\psi})^{N}
    + (N-1)\left(1 - \frac{F(v_{\psi})}{F(z)}\right) F(z)^{N}
    - F(v_{\psi})\bigl(F(z)^{N-1} - F(v_{\psi})^{N-1}\bigr)
  \notag\\
  &= N\bigl(F(z) - F(v_{\psi})\bigr) F(z)^{N-1},
  \label{eq:profit-fail-bracket}
\end{align}
where the second line uses $F(v_{\psi})^{N}=F(v_{\psi})F(v_{\psi})^{N-1}$.
When $N=1$, only the first integral enters and
\eqref{eq:profit-fail-bracket} still holds.
Substituting into \eqref{eq:profit-fail},
\begin{equation}
  \Pi_{\text{fail}}
  = \delta N (c_H - c_L)\,
    \bigl(F(z) - F(v_{\psi})\bigr)\,
    F(z)^{N-1}.
  \label{eq:profit-fail-closed}
\end{equation}
Raising $\underline{r}$ makes stage~1 failure more likely because $F(z)^{N}$
increases, and for $\underline{r}>\underline{r}_0$ it also enlarges
$[v_{\psi},z]$ and raises the non-max offer probability
$1-F(v_{\psi})/F(v)$.

\subsection{When the auction succeeds}
\label{subsec:profit-success}

The auction \emph{succeeds} when the reserve is met in stage~1, so the sale
does not move to stage~2.
This is equivalent to $v_{\max}>z$, with probability $1-F(z)^{N}$.
Once the reserve is met, bidders bid truthfully in the English auction.
The highest-value bidder $v$ wins and pays $\max\{\underline{r},s\}$, so the
seller's margin is $\max\{\underline{r}, s\} - c_L$.
Averaging over the winner's value $v\ge z$ and the runner-up value $s\le v$ gives
\begin{equation}
  \Pi_{\text{succ}}
  = \int_{z}^{1}
    \int_{c_L}^{v}
      \bigl(\max\{\underline{r}, s\} - c_L\bigr)
      \, d\!\left(\frac{F(s)}{F(v)}\right)^{N-1}
    \, dF(v)^{N},
  \label{eq:profit-success-int}
\end{equation}
where $d(F(s)/F(v))^{N-1}$ is the distribution of $s$ conditional on
$v_{\max}=v$.
Integrating by parts yields the equivalent form
\begin{equation}
  \Pi_{\text{succ}}
  = N \Biggl\{
    (\underline{r} - c_L) F(\underline{r})^{N-1}
      \bigl(1 - F(z)\bigr)
    + \int_{z}^{1}
      \int_{\underline{r}}^{v} (s - c_L)\, dF(s)^{N-1}
      \, dF(v)
  \Biggr\}.
  \label{eq:profit-success}
\end{equation}
The first term collects states where the runner-up's value is below
$\underline{r}$, so the sale price equals the reserve. The double integral
collects states where $s>\underline{r}$, so the winner pays the runner-up's
value $s$ (the usual second-price payment in an English auction).

\subsection{Total profit and \texorpdfstring{$\underline{r}^{*}$}{underline r star}}
\label{subsec:total-profit}

Total profit is
\begin{equation}
  \Pi(\underline{r})
  = \mathbf{1}\{\underline{r}>\underline{r}_0\}\,
    \delta N (c_H - c_L)\,
    \bigl(F(z) - F(v_{\psi})\bigr)\,
    F(z)^{N-1}
    + \Pi_{\text{succ}},
  \label{eq:Pi-low-total}
\end{equation}
The indicator $\mathbf{1}\{\underline{r}>\underline{r}_0\}$ makes explicit the
reserve cutoff below which stage-2 profit is zero.
Stage~2 offers always use the buyers' posterior belief $\tilde{\psi}$.
In \eqref{eq:Pi-low-total}, failure before $\underline{r}$ leaves
$\tilde{\psi}=\psi$, so the cutoff is $v_{\psi}$.
If the seller deviates upward, stage-2 payoffs use the updated posterior
$\tilde{\psi}$ from \eqref{eq:belief-update}, which is strictly below $\psi$.

\begin{claim}
  \label{claim:underline-r-star}
  The equilibrium lower support $\underline{r}^{*}$ is the argmax of
  $\Pi(\cdot)$ in \eqref{eq:Pi-low-total}.
\end{claim}

\subsection{Proof of Claim~\ref{claim:underline-r-star}}
\label{subsec:proof-r-star}

\begin{proof}
Let $\underline{r}^{*}\in\arg\max\Pi(\cdot)$.  We show that any other lower
support $\underline{r}\neq\underline{r}^{*}$ cannot be optimal.

\smallskip\noindent
\textit{Downward deviation.}
If $\underline{r}>\underline{r}^{*}$, failure still occurs before $\underline{r}$,
so the posterior remains $\tilde{\psi}=\psi$ and
\eqref{eq:Pi-low-total} is exact, so $\Pi(\underline{r}^{*})>\Pi(\underline{r})$.

\smallskip\noindent
\textit{Upward deviation.}
If $\underline{r}<\underline{r}^{*}$, announce $\underline{r}^{*}$ and write
$z\equiv\beta^{-1}(\underline{r})$ and $z^{*}\equiv\beta^{-1}(\underline{r}^{*})$.

\smallskip\noindent
\emph{Case 1: $\underline{r}^{*}>\underline{r}>\underline{r}_0$.}
Then $z>v_{\psi}$ and $z^{*}>z$.
True failure profit is
\begin{align}
  \Pi_{\text{fail}}^{\text{true}}
  &= \mathbf{1}\{\underline{r}>\underline{r}_0\}\,
    \delta N (c_H - c_L)\,
    \bigl(F(z) - F(v_{\psi})\bigr)\,
    F(z)^{N-1}
  \notag\\
  &\quad + \delta (c_H - c_L)
    \int_{z}^{z^{*}} dF(v)^{N}
    + \delta (N-1) (c_H - c_L)
    \int_{z}^{z^{*}}
      \left(1 - \frac{F\bigl(v_{\rho}(v)\bigr)}{F(v)}\right)\, dF(v)^{N},
  \label{eq:Pi-true-fail}
\end{align}
where the first term is $\Pi_{\text{fail}}$ and
$v_{\rho}(v)$ is the $c_H$-offer cutoff with belief
$\tilde{\psi}(\beta(v))$ on $(z,z^{*}]$.
On $(z,z^{*}]$, $\tilde{\psi}(\beta(v))<\psi$, so
$v_{\rho}(v)<v_{\psi}<z<z^{*}$ and the last two terms are strictly positive.
The failure term in $\Pi(\underline{r}^{*})$ uses $v_{\psi}$ up to $z^{*}$, so
the updated-belief block strictly exceeds its contribution.
Success profit at $\underline{r}^{*}$ is unchanged, hence
$\Pi_{\text{fail}}^{\text{true}}>\Pi(\underline{r}^{*})>\Pi(\underline{r})$.

\smallskip\noindent
\emph{Case 2: $\underline{r}^{*}>\underline{r}$ but $\underline{r}\le\underline{r}_0$.}
Then $\mathbf{1}\{\underline{r}>\underline{r}_0\}=0$ and $\Pi_{\text{fail}}=0$.
If $\underline{r}^{*}>\underline{r}_0$, the integrals over $(z,z^{*}]$ in
\eqref{eq:Pi-true-fail} are positive; otherwise compare success profit only.

\smallskip\noindent
\textit{Conclusion.}
Either deviation is profitable, so $\underline{r}^{*}$ is the optimal lower
support.
\end{proof}

The low-cost seller's equilibrium payoff is $\Pi(\underline{r}^{*})$. On the
low branch,
$\beta^{-1}(\underline{r})=(\underline{r}-\delta\psi c_L)/(1-\delta\psi)$.
